\documentclass[../main.tex]{subfiles}

\graphicspath{{\subfix{../images/}}}



\begin{document}

%% BEGIN SUBFILE %%

\newpage

\section{Electronics measures}

\subsection{Introduction}

In this first section, we will introduce the main concepts related to measurements in general, and then apply them to the field of electronics, which is our area of specialization. In particular, we will analyze measurements and their uncertainties in relation to possible requirements; we will then present the measuring instruments that will be used in the laboratory.\\

Let us now introduce the concept of measurement and, more generally, the reason why it is necessary to measure certain quantities, especially in engineering or science.

Let's imagine we need to develop a product with well-defined specifications and requirements to meet. Very often, these requirements are expressed in the form of measurements (e.g., dimensions, performance, consumption, tolerances). Furthermore, considering a more complex context in which competition also exists, it becomes essential to establish which product is best. In order to make an objective judgment through comparison, it is essential to take measurements with adequate characteristics in terms of accuracy, precision, and reliability.

We can therefore say that, even in the case of a single product or any system whose performance over time we want to evaluate - in terms of improvement or deterioration - measurement is an essential tool. Ultimately, it is not possible to determine an improvement (or deterioration) without the ability to measure correctly and consistently.\\

With this little introduction let's now dig a bit deeper inside the history of measuring things and how it evolved. Historically, humanity defined units using pieces of metal, the size of a king's foot, or the rotation of a wobbly planet. But on May 20, 2019, metrology (the science of measurement) underwent the greatest revolution since the French Revolution birthed the metric system.

We threw away the physical artifacts. We stopped defining units by things and started defining them by the fundamental constants of the universe. Let us dissect this magnificent system piece by piece.

\subsection{The paradigm shift, from artifacts to constants}

Before 2019, the definition of units was backwards. We had a platinum-iridium cylinder in a vault in Paris called the International Prototype of the Kilogram (IPK), or "Le Grand K." By definition, its mass was exactly 1 kilogram. The problem? If someone dropped Le Grand K, or if it absorbed a microgram of contamination, the mass of the entire universe (measured in kg) mathematically shifted. The standard was vulnerable, local, and degrading.

The brilliant conceptual leap of the modern SI system is this: We took seven fundamental constants of nature, gave them exact numerical values with zero uncertainty, and forced the units to adapt to the constants.

If we ever meet an advanced alien civilization, they will not know what Le Grand K is. But they will know the speed of light, and they will know the charge of an electron. The modern SI system is a language written for the cosmos.

\subsection{The seven base units}

Every single measurement in the universe—from the torque of a motor to the impedance of a trace on a PCB, to the luminosity of a star—is derived from just seven base units. Let us explore where each one comes from.

\paragraph{Time: second (s)}

\begin{enumerate}
	\item[\textbf{-}] The Old Way: 1/86400th of a mean solar day. Problem: The Earth’s rotation is slowing down due to tidal friction and wobbles unpredictably. A terrible clock. The Modern Definition: Tied to the hyperfine transition frequency of the Cesium-133 atom ($\Delta\nu_{Cs}$).
	\item[\textbf{-}] How it works: in quantum mechanics, electrons exist in discrete energy levels. If you excite a Cesium-133 atom, its outermost electron jumps to a higher state, then falls back, emitting a photon of microwave radiation. Because atomic structure is identical across the universe, this frequency never changes.
	\item[\textbf{-}] The Exact Rule: we defined the unperturbed ground-state hyperfine transition frequency of the Cesium-133 atom to be exactly 9,192,631,770 Hz (cycles per second).
	\item[\textbf{-}] The Realization: we build "atomic fountains" where lasers cool and toss Cesium atoms upward, bombarding them with microwaves until they resonate. We simply count 9'192'631'770 cycles, and we call that exactly one second.
\end{enumerate}

\paragraph{Length: meter (m)}

\begin{enumerate}
	\item[\textbf{-}] The Old Way: a fraction of the Earth's meridian, then a platinum-iridium bar, then the wavelength of Krypton-86 gas.
	\item[\textbf{-}] The Modern Definition: tied to the speed of light in a vacuum ($c$).
	\item[\textbf{-}] How it works: thanks to Albert Einstein, we know the speed of light in a vacuum is the absolute speed limit of the universe. It is a perfect constant.
	\item[\textbf{-}] The Exact Rule: we defined the speed of light to be exactly 299'792'458 meters per second.
	\item[\textbf{-}] The Realization: since we already defined the second (via Cesium), the meter is simply the distance light travels in a vacuum in exactly 1/299'792'458 of a second.
\end{enumerate}

\paragraph{Mass: kilogram (kg)}

\begin{enumerate}
	\item[\textbf{-}] The Old Way: the physical chunk of metal in Paris (Le Grand K).
	\item[\textbf{-}] The Modern Definition: tied to the Planck constant ($k$).
	\item[\textbf{-}] How it works: this was the hardest unit to fix. How do you tie mass to a quantum mechanical constant? You use Albert Einstein ($E = mc^2$) and Max Planck ($E = hf$). Mass is equivalent to energy, and energy is proportional to frequency via the Planck constant.
	\item[\textbf{-}] The Exact Rule: we defined the Planck constant $h$ to be exactly $6.62607015 \cdot 10^{-34}$  Joule-seconds (which is $\unit{kg} \cdot \unit{m^2 / s} $).
	\item[\textbf{-}] The Realization (The Kibble Balance): this is a masterpiece of engineering. A Kibble balance weighs a mass not against another mass, but against an electromagnetic force. We push a coil of wire through a magnetic field, creating a voltage (measured flawlessly using the quantum Josephson effect). We run a current through that coil to hold the mass up against gravity (current is measured using the quantum Hall effect). By balancing mechanical power ($mgv$) with electrical power (VI), the mass mathematically resolves directly to the Planck constant. No physical weights required.

\end{enumerate}

\paragraph{Electric Current: Ampere (A)}

\begin{enumerate}
	\item[\textbf{-}] The Old Way: the current that, if maintained in two straight parallel wires of infinite length, 1 meter apart in a vacuum, produces a force of $2 \cdot 10^{-7}$ Newtons per meter. Problem: Infinite wires do not exist. It was an un-testable nightmare for actual metrologists.
	\item[\textbf{-}] The Modern Definition: tied to the elementary charge ($e$), the charge of a single proton (or absolute value of an electron).
	\item[\textbf{-}] How it works: current is just the flow of electric charge over time ($I = dQ / dt$).
	\item[\textbf{-}] The Exact Rule: we defined the elementary charge to be exactly $1.602176634 \cdot 10^{-19}$ Coulombs (Ampere-seconds).
	\item[\textbf{-}] The Realization: we use Single-Electron Pumps (SEPs). These nano-devices trap individual electrons in quantum dots and push them through a circuit one by one at a known frequency (using our Cesium clock). If you pump exactly $6.241509 \cdot 10^{18} $electrons per second, you have exactly one Ampere.
\end{enumerate}

\paragraph{Thermodynamic temperature: Kelvin (K)}

\begin{enumerate}
	\item[\textbf{-}] The Old Way: 1/273.16 of the thermodynamic temperature of the triple point of water. Problem: water isotopic composition varies depending on where on Earth you get it.
	\item[\textbf{-}] The Modern Definition: tied to the Boltzmann constant ($k_B$).
	\item[\textbf{-}] How it works: temperature is not a fundamental property; it is a macroscopic illusion. What you feel as "heat" is actually the average kinetic energy of microscopic particles vibrating and flying around. The Boltzmann constant bridges the microscopic energy of particles with the macroscopic concept of temperature.
	\item[\textbf{-}] The Exact Rule: we defined $k_b$ to be exactly $1.380649 \cdot 10^{-23}$ Joules per Kelvin.
	\item[\textbf{-}] The Realization: we measure temperature by measuring the acoustic resonance of gas inside a perfectly spherical copper cavity (Acoustic Gas Thermometry), or by measuring the thermal electrical noise in a resistor (Johnson-Nyquist noise). We measure the pure energy, apply $k_B$, and out pops the temperature in Kelvin.
\end{enumerate}


\paragraph{Amount of substance: mole (mol)}

\begin{enumerate}
	\item[\textbf{-}] The Old Way: the number of atoms in exactly 12 grams of Carbon-12. Problem: This tied chemistry to the definition of the kilogram, which was tied to the metal cylinder in Paris.
	\item[\textbf{-}] The Modern Definition: tied to the Avogadro constant ($N_A$).
	\item[\textbf{-}] How it works: a mole is literally just a number, like a "dozen" (12). It is used by chemists to translate between the mass of a substance and the number of atoms in it.
	\item[\textbf{-}] The Exact Rule: we defined $N_A$ to be exact to be exactly $6.02214076 \cdot 10^{23}$ entities per mole.
	\item[\textbf{-}] The Realization: it is now purely a counting number. If you have exactly $6.02214076 \cdot 10^{23}$ electrons, you have one mole of electrons. The link to Carbon-12 is severed; the mole stands alone as a pure numeric constant.
\end{enumerate}

\paragraph{Luminous intensity: candela (cd)}

\begin{enumerate}
	\item[\textbf{-}] The Old Way: the brightness of a specific type of candle, then the brightness of freezing platinum.
	\item[\textbf{-}] The Modern Definition: tied to the luminous efficacy of monochromatic radiation of frequency $540 \cdot 10^{12}$ Hz ($K_{cd}$).
	\item[\textbf{-}] How it works: this is the odd one out. It is the only unit tied to human biology. The frequency 540 THz is yellowish-green light, the exact color to which the human eye is most sensitive. The Candela bridges objective radiometric power (Watts) with subjective human visual perception.
	\item[\textbf{-}] The Exact Rule: $K_{cd}$ is exactly 683 lumens per Watt at that specific frequency.
	\item[\textbf{-}] The Realization: we use radiometers to measure absolute optical power in Watts, and then apply the defined human-eye sensitivity curve to yield Candelas.
\end{enumerate}


\subsubsection{Derived units} % TODO - Small work


\subsection{Measures}

\subsubsection{Uncertainty} % TODO MAYBE REMOVE SUB SUB SECTION?


\subsection{Oscilloscope}

After providing a comprehensive introduction to the concept of measurement and the associated uncertainties, we can now proceed to analyze one of the main measuring instruments used in laboratories. This instrument allows us to take measurements and observe what happens inside a circuit that we want to analyze.

\subsubsection{Basic operation}

Let's now begin to analyze the way an oscilloscope work starting by looking at how an analog one works. From seeing how the analog version works we will then move on to analyze and study the property and functions of its digital counterpart.


\subsubsection{Display}

\subsubsection{Examples}

\subsection{Digital sampling oscilloscope - DSO}

\subsubsection{Quantization}


\subsubsection{Sampling}


\subsubsection{Oscilloscope connections}

After introducing the operation of the oscilloscope, both from a theoretical and practical point of view, we can now analyze how it is connected to the outside world and what operations need to be performed before taking a measurement.\\

In particolare l'oscilloscopio digitale a campionamento utilizza dei connetttori denominati BNC (Bayonet Neill–Concelman) utilizzati principalmente per quick-connect/disconnect, locking connector used for coaxial cables featuring a bayonet-style, twist-lock mechanism that prevents accidental disconnection, making it ideal for high-frequency. 

At this point we might imagine that being the probe some kind di circuiti se stessa deve essere tenuta in considerazione all'interno della nostra misura in quanto potrebbe agire in un determianto modo sul fenomeno che andiao a misurare.

In particoalre, essendo nel modno reale, il connettore in ingresso non è ideale: essa infatta introduce una resistenza che denominiam ocome $R_i$ e una capacità denominata come $C_i$. Possiamo dunque modellare il connettore come una resistenza in parallelo ad un condensatore come riportato di seguito.

% TODO MODELLO R + C

Arrivati a questo punto possiamo agigungere al nostro modello un generatore di tensione messo ad indicare la funzione che viene posta in ingresso al nostro oscilosscopio. Dato che anche questo non è ideale è necessario aggoungere un'opportuna resistenza $R_g$ in serie al generatore stesso. Infine è doveroso introdurre, in quanto viene utilizzato un cavo per connetere questi due componenti, una capacità parassita aggiunta proprio da questa ultimo componente. Mettendo le varie parti assieme otteniamo il seguente schema:

% TODO SCHEMA COMPLETO 1

In aggiunta a quanto detto possiamo andare a dare degli ordini di grandezza della resistenza $R_i$ e della capacità $C_i$ che si aggirano attorno ad 1 M$\Omega$ ed a decine di pF.\\

Sapendo l'ordine di grandezza di questi due parametri possiamo avere un'idea del valore visto e misurato dall'oscilloscopio ai capi della capacità $C_i$.
Analizzando il circuito visto sopra arriviamo alla seguente espressione:

\[ V_i = V_g \cdot \frac{R_i}{R_i + R_g} \cdot \frac{R_i}{s(R_g \;//\; R_i)C_p + 1}\]\

Arrivati a questo punto possiamo fare l'assunzione che il valore di $R_g$, essendo una resistenza parassita, sia molto bassa e duqnue molto minore di $R_i$: $R_g \ll R_i$.
Data questa condizione possiamo allora andare a dire che il valore dell'attenuazione data dal partitore di tensione (primo termine moltiplicativo) sia approssimabile ad uno.
Segue quanto detto in termnimi matematici:

\[ R_g \ll R_i \Rightarrow \frac{R_i}{R_i + R_g}  \simeq 1 \]\

Arrivati a questo punto, fatta la precedente assunzione sui valori di $R_g$ e $R_i$ possiamo dire che in corrente dirett (DC) il valore misurato a capi di $C_i$, ovvero $V_i$, sia circa pari a quello in ingresso: $V_g \simeq V_i$.\\

Sapendo che però, come visto negl iesempi precedneti, sappiamo che con lo struento andiao anda alizzare forem d'onta di natura periodica ci chiadiamo come questo modello di circuito si comporti nel caso d itensioni alternate nel tempo.

Possiamo dunqiue fare l'analisi del polo del secondo membro moltiplicato si ottiene che:

\[ f_t = \frac{1}{2\pi (R_g \;//\; R_i)C_p}\]\

Operando sempre attraverspo la semplificazione fatta prima otteniamo che:

\[ f_t = \frac{1}{2\pi  R_i C_p}\]\

Inserendo ora dei possibili valori, di cui abbimo visto prima i possiibli ordini di grandezza, possiamo andare ad avere un'idea della frequenza di taglio.
Ipotizzando che $R_g$ sia 50 $\Omega$ e che $C_p$ sia 100 pF si ottiene il seguente risultato:

\[ f_t = \frac{1}{2\pi  R_i C_p} = \frac{1}{2\pi \cdot 50 \Omega \cdot 100 \unit{pF}} = 31.830\; \unit{ks^{-1}} \simeq 32 \unit{MHz} \]\


Come possiamo vedere abbiamo raggiunto, per le proiprità fisiche del nostro struemnto, un limite che non possiamo sorpassare. In altre parole non siamo in grado con il setup corrente di andare a misurare senza avere attenuazioni significative segnale con frequenze superiori ad i 32 MHz.\\

Per ovviare a questo problea possiamo andare ad intordure una seconda componente all;intenro del nostro circuito per costruire una sonda compensata, ovvero che compensi l'effetto ch abbiamo appena sottolineato.

Al fine di realizzare una sonda compensata andiamo ad aggiungere un dipolo compsoto da una resistenza $R_s$ connessa in parallea ad una capacotà variabile $C_s$ come mostrato di seguito:

% TODO CIRCUITO COMPENSATO

A quesot punto possiamo andare a calcolare la tensione ottenuta ai capi del consendatore del connettore. Si ottiene che:

\[ V_i = V_g \cdot \frac{\dfrac{R_i}{sR_iC_p + 1}}{R_g + \dfrac{R_s}{sR_sC_s + 1} + \dfrac{R_i}{sR_iC_p + 1}}\]\

Come visto in precedenza andiamo a vedere come si comporta questo circuito con correnti dirette e non. Prima di tutto andiamo a fare delle assunzioni:

\begin{enumerate}
	\item[\textbf{-}] $C_p = C_c + C_i $
	\item[\textbf{-}] $R_s \cdot C_s = R_i \cdot C_i $
	\item[\textbf{-}] $R_s = 9C_p$ e dunque $C_s = \dfrac{C_p}{9}$	
	
\end{enumerate}

Fatte queste premesse edassunzioni arriviamo a dire che:

\[ V_i = V_i \dfrac{R_i}{R_g + 10R_i}  \cdot \dfrac{1}{\dfrac{s R_gR_i}{R_g + 10R_i}C_p + 1} \]\


A questo punto in corrente diretta (DC) si ha che:

\[ \dfrac{R_i}{R_g + 10R_i}  \simeq \dfrac{R_i}{10R_i} \simeq \frac{1}{10} \]\

Analizzando invece il circuito a tensioni non costanti si ottiene che:

\[ f'_t = \frac{1}{2\pi \dfrac{R_gR_i}{R_g + 10R_i}C_p } \simeq \frac{1}{2\pi \dfrac{R_g}{10}C_p }  = \frac{10}{2\pi R_g C_p} \]\

Possiamo notare come calcolando il valore ottenuto di $f'_t$ si ottiene una frequenza dieci volte più grande e dunque una maggiore banda di segnale in ingresso che non viene attenuata.
Svolgendo i calcoli si ottiene che:


\[ f'_t = \frac{10}{2\pi R_g C_p} = \frac{10}{2\pi \cdot 50 \Omega \cdot 100 \unit{pF}} = 318.30\; \unit{ks^{-1}} \simeq 320 \unit{MHz} \]\

A questo punto comopreso andiamo a vedere nella pratica cosa succede. In partcioalre abbiamo visot come una sonda compensta carichi di meno il circuito e che quindi influisca con un'attenuazione costante. Infine dato che la capacità del compensatore deve essere variabile in quanto il suo valore di compensazione diopend dal cavo utilizzato e dalla capacità in ingresso dell'oscilloscopio. Viene dunque introdotto, cme visot in figura, un capacità variabile direttamente dall'utente affichè funzioni corettamnete.

Andnado ne lparticoalr la regolazione avviene con un cacciavite in plastica - uno in metallo aggiungerebbe l'ennesima capacità parassita - sulla sonda stessa mentre essa è collegata a terra ed ad un'apposita uscita sull'oscilloscopio (PROBE COMP) che genera un'onda quadra avente certe caratteristiche tipo: $V_{pp}$ 1-5 V con frequenza $f_c$ pari a 1 kHz.
Dato che se la sonda non è ben compensata si comporta come un filtro passa basso l'onda quadra generata verrà visualizzata sullo schemro dell'osciloscoapio affeta da alcuni difetti sui fronti di salita e discesa. È ragionevole pensare ciò in quanto queste parti del segnale rappresentano nel dominio della frequenza le componenti a maggiore frequenza. Di seguoito sono riportati degli esempi di onde quadre sovra e sotto-compensate. 

% TODO IMAGE SOV E SOTT COMPENSAZIONE

Notiamo come la sottcompensazione porti ad un drift crescente della parte piatta dell;onda  metre la sovracomoensazione porti la genereazione di un picco dopo  

Tornando a noi una volta visualizzata l'onda sullo strumento ci è possibile regolare la capacità fino a quando non visualizzaimo una corretta onda quadra in che significa che la sonda è compensata correttamente.


\subsubsection{Examples}


\subsection{Digital Multimeter - DMM}


\subsubsection{Examples}


\begin{center}
\begin{tikzpicture}[transform shape]
	% Paths, nodes and wires:
	\draw (12.5, 5) to[generic, mirror, v=$V_{out}$, voltage/distance from node=0.4, voltage/shift=2] (12.5, 8);
	\draw (6, 8) to[capacitor, l={$C_{in}$}] (6, 5);
	\node[nmos, rotate=-90](N1) at (6.98, 8){} node[anchor=north] at (N1.text){$SW_1$};
	\draw (8, 5) to[empty diode, l_={$SW_2$}] (8, 8);
	\draw (8, 8) to[cute inductor, l={$L$}, v<=$V_L$, voltage/shift=1] (11, 8);
	\draw (5, 8) -- (6.25, 8);
	\draw (7.75, 8) -- (8, 8);
	\draw (11, 8) -- (12, 8);
	\draw (12, 8) -- (12.5, 8);
	\draw (12.5, 5) -- (12, 5);
	\draw (11, 8) to[capacitor, l_={$C_{out}$}, label distance=0.07cm] (11, 5);
	\draw (12, 5) -- (6, 5);
	\draw (5, 8) to[open, v_<=$V_{in}$, voltage/bump b=1, voltage/shift=-0.4] (5, 5);
	\node[ocirc] at (5, 8){};
	\node[ocirc] at (5, 5){};
	\node[ocirc] at (6.98, 8.98){};
	\draw (5, 5) -- (6, 5);
\end{tikzpicture}
\end{center}



































\end{document}



